% \begin{savequote}
% \includegraphics[width=3.5cm]{./images/pallinger-christoph-aass/Anton_32986_v2-AASS-0112mm.jpg}
% 
% Bild: Christoph Pallinger
% \end{savequote}
\def\ChapterCode{ASS}
\chapter
[\CO{[ASS]} Assistenztätigkeiten]
{Assistenztätigkeiten}
% \AasRsN{RS~XIII}
\label{chp:ASS}

\ChapterIntro
{%
~~
}
{%
\begin{description}
  \item[Maintainer:] Diverse
  \item[Autoren:] Diverse
  \item[Reviewer:] Standard-Reviewprozess
  \item[Version:] {\VarDocVersionTypeName} ({\VarDocVersionTypeDescription})
  \item[SHA1:] {\DocGitCommit}
\end{description}

{\TxTeilkapitelAASS}
}


\Text{Special Credits}{%
\noindent Teile des Textes und der Handlungsanleitungen sind
in Zusammenarbeit mit dem \emph{Department für medizinische
Aus- und Weiterbildung -- Medizinische Universität Wien,
Abt. Methodik und Entwicklung},  welches in Zusammenarbeit 
mit dem Klinischen Institut für Krankenhaushygiene
entstanden. Wir danken Ass.-Prof. Dr. Michael Schmidts, Dr.
Alexander Blacky und deren Team für die Unterstützung!
}{

}
{icons/muw.png}
{}
{MUW}





\clearpage
\section{Arbeiten mit  \HeadingKeyword{Medikamenten}}

\subsection{\HeadingKeyword{Allgemeines} zu Medikamenten}

%
%
%\T{Arzneimittel} sind Stoffe oder Zubereitungen aus Stoffen, 
%die zur Anwendung im oder am Körper 
%und als Mittel mit Eigenschaften zur Heilung, Linderung oder Verhütung 
%von Krankheiten oder krankhafter Beschwerden bestimmt sind, 
%bzw. die physiologische Funktionen durch eine pharmakologische, immunologische oder metabolische Wirkung wiederherstellen, korrigieren oder beeinflussen können, 
%oder als Grundlage für eine medizinische Diagnose zu dienen.
%Sie enthalten \T{Wirkstoffe}, das sind Stoffe oder Gemische von Stoffen, die die arzneilich \E{wirksamen Bestandteile des Arzneimittels} sind.
%Daneben können sie \T{Hilfsstoffe} enthalten, das sind alle Bestandteile eines Arzneimittels mit Ausnahme der Wirkstoffe und des Verpackungsmaterials.
%
%\T{Arzneispezialitäten} sind Arzneimittel, die im Voraus stets in gleicher Zusammensetzung hergestellt und unter der gleichen Bezeichnung in einer zur Abgabe an den Verbraucher oder Anwender bestimmten Form in Verkehr gebracht werden.
%Der \T{Name des Arzneimittels} ist der Name, der entweder ein nicht zu Verwechslungen mit dem gebräuchlichen Namen führender \E{Phantasiename} oder ein gebräuchlicher oder \E{wissenschaftlicher Name} in Verbindung mit einem Warenzeichen oder dem Namen des Zulassungsinhabers sein kann.
%Ein \T{Generikum} ist ein Arzneimittel, das die gleiche qualitative und quantitative Zusammensetzung aus Wirkstoffen und die gleiche Darreichungsform wie das Referenzarzneimittel aufweist und dessen Gleichartigkeit mit dem Referenzarzneimittel durch geeignete Studien nachgewiesen wurde\ZFootnote{Bioäquivalenz}.  
%
%Arzneimittel haben eine begrenzte \T{Haltbarkeit}, das ist die Eigenschaft, während eines bestimmten Zeitraumes bei ordnungsgemäßer Lagerung die Beschaffenheit, insbesondere im Hinblick auf Qualität und Wirkung, nicht zu verändern. Das \T{Verfalldatum} ist die Bezeichnung jenes Zeitpunktes, nach dem die Haltbarkeit eines Arzneimittels nicht mehr gewährleistet ist.
%
%„Nebenwirkung“ eines Humanarzneimittels ist eine schädliche und unbeabsichtigte Reaktion auf das Arzneimittel. „Nebenwirkung“ eines Tierarzneimittels ist eine schädliche und unbeabsichtigte Reaktion auf das Arzneimittel, die bei Dosierungen auftritt, wie sie normalerweise bei Tieren zur Prophylaxe, Diagnose oder Therapie von Krankheiten oder für die Änderung einer physiologischen Funktion verwendet werden.
%
%Das Bundesamt für Sicherheit im Gesundheitswesen führt ein Register (\I{Arzneispezialitätenregister}) in welches zugelassene Arzneispezialitäten eingetragen werden und in dem die zugehörigen Fachinformationen abrufbar sind. Es ist im Internet zugänglich unter der URL \url{https://aspregister.basg.gv.at/}.
%
%
%
%Pflichten von Angehörigen der Gesundheitsberufe
%
% 75g. (1) Ärzte, Zahnärzte, Tierärzte, Dentisten, Hebammen und, soweit sie nicht der Meldepflicht gemäß  75j unterliegen, Apotheker und Gewerbetreibende, die gemäß der Gewerbeordnung 1994 zur Herstellung von Arzneimitteln oder zum Großhandel mit Arzneimitteln berechtigt sind, und Drogisten haben
%										
%
%1.
%	
%
%vermutete Nebenwirkungen oder
%
%2.
%	
%
%vermutete Nebenwirkungen beim Menschen oder
%
%3.
%	
%
%das Ausbleiben der erwarteten Wirksamkeit oder
%
%4.
%	
%
%nicht ausreichende Wartezeiten
%
%von Arzneimitteln, die im Inland aufgetreten sind und ihnen auf Grund ihrer beruflichen Tätigkeit bekannt geworden sind, nach Maßgabe einer Verordnung gemäß  75a unverzüglich dem Bundesamt für Sicherheit im Gesundheitswesen zu melden.
%
%(2) Die gemäß Abs. 1 Meldepflichtigen haben dem Bundesamt für Sicherheit im Gesundheitswesen alle Beobachtungen und Daten mitzuteilen, die für die Arzneimittelsicherheit von Bedeutung sein können.
% 75h
%Text
%
% 75h. Personen, die nicht der Meldepflicht nach  75g und 75j unterliegen, insbesondere Patienten, haben die Möglichkeit, vermutete Nebenwirkungen von Arzneimitteln dem Bundesamt für Sicherheit im Gesundheitswesen über das Internetportal für Arzneimittel oder auf dem Postweg zu melden.
%
%Verstöße gegen Abgabebedingungen o.\,ä. des Arzneimittelgesetzes sind mit hohen Geldstrafen sanktionier (gegenwärtig bis zu 7\,500 Euro, im Wiederholungsfalle bis zu 14\,000 Euro).


\Text
{Stärke: Wirkstoff, Dosis, Konzentration}
{Die \T{Stärke des Arzneimittels} ist je nach Verabreichungsform der Wirkstoffanteil pro Dosierungs-, Volumens- oder Gewichtseinheit.
%
Bei Medikamenten wird normalerweise die Menge des
Wirkstoffes (z.\,B. in \Mg) angegeben. 
%

Manchmal erfolgt die
Angabe auch in anderen Einheiten, wie z.\,B. den
\I{Internationalen Einheiten} (\I{IE}). Bei Ampullen findet sich außerdem auch eine Volumsangabe.
Aus der Angabe der Wirkstoffmenge und des Volumens kann man
die Konzentration berechnen.
%

In manchen Fällen ist auch eine Konzentrationsangabe üblich 
(z.\,B. 
Glukose \SI{5}{\percent}, 
NaCl \SI{0,9}{\percent}, 
Propofol \SI{1}{\percent}, Adrenalin \Conc{1}{1000} 
usw.). 
%
Ein Prozent enspricht dann umgerechnet \SI{10}{\milli\gram\per\milli\liter}
bzw. \SI{1}{\milli\liter} entspricht \SI{1000}{\milli\gram}.


\TableWide
{Typische Angaben zu Wirkstoffmenge und Konzentration}%1 Titel
{Beachte: 
Gleichlautende Spezialitäten können erheblich unterschiedliche Wirkstoffkonzentrationen sowie Wirkstoff- und Füllmengen aufweisen!}%2 Beschreibung
{}%3 Label/Hook
{}%4 Quelle
{}%5 Lizenz
{%
\begin{tabulary}{\linewidth}{LRRCL}
\toprule\addlinespace
\noindent\TabHeading{Form} 	& \noindent\TabHeading{Wirkstoffmenge} 	& \noindent\TabHeading{Volumen} 	& \noindent\TabHeading{Konzentration}							& \noindent\TabHeading{Beispiele}
\\\addlinespace\midrule\addlinespace
Ampulle 			& 2\Mg* 						& 20\Ml* 				& \SI{0,1}{\mg\per\ml} \newline=\,\SI{1}{\mg\per{10\,}\ml}
\newline =\,\Conc{1}{10\,000}	& L-Adrenalin
\\\addlinespace
Ampulle 			& 5\Mg* 						& 5\Ml* 				& \SI{1}{\mg\per\ml}								& Dormicum{\texttrademark}
\\
Ampulle 			& 15\Mg* 						& 3\Ml* 				& \SI{5}{\mg\per\ml}								& Dormicum{\texttrademark}
\\\addlinespace
Ampulle 			& 25\Mg* 						& 5\Ml* 				& \SI{5}{\mg\per\ml}								& Ketanest{\texttrademark}
\\
Ampulle 			& 50\Mg* 						& 2\Ml* 				& \SI{25}{\mg\per\ml}								& Ketanest{\texttrademark}
\\\addlinespace
Stechampulle 			& 500\Mg* 						& 50\Ml* 				& \SI{10}{\mg\per\ml} 
=\,\SI{1}{\percent}	& Propofol
\\
Stechampulle 			& 1\,000\Mg* 						& 50\Ml* 				& \SI{20}{\mg\per\ml} 
=\,\SI{2}{\percent}	& Propofol
\\\addlinespace
Tabletten 			& 500\Mg* 						&						&														& Parkemed{\texttrademark}, Mexalen{\texttrademark}, \ldots
\\\addlinespace
Trockenstechampulle & 4000 IE						&						&														& Heparin
\\\addlinespace\bottomrule\addlinespace
\end{tabulary}
}%6 Tabelle
{}%7 Float


% \begin{itemize}
%   \item Ampulle mit Wirkstoff 2\Mg* in 20\Ml* {\dotfill}
%   Konzentration 0,1\nicefrac{\Mg}{\Ml} oder
%   1\nicefrac{\Mg}{10\Ml*}
%   \item Ampulle mit Wirkstoff 5\Mg* in 5\Ml* {\dotfill}
%   Konzentration 1\,\nicefrac{\Mg}{\Ml}
%   \item Ampulle mit Wirkstoff 15\Mg* in 3\Ml* {\dotfill}
%   Konzentration 5\,\nicefrac{\Mg}{\Ml}
%   \item Tabletten mit Wirkstoff 500\Mg* (pro Tablette)
%   \item Trockenstechampulle mit Wirkstoff 4000 IE
% \end{itemize}

Viele Medikamente können auch zur besseren Dosierung
mit \EE{geeigneten Substanzen} (siehe
Packungsbeilage) verdünnt werden:
\begin{itemize}
  \item Ampulle mit Wirkstoff 15\Mg* in 3\Ml* 
  {\dotfill} Konzentration 5\,\nicefrac{\Mg}{\Ml}
  \item Obige Ampulle 
  + 12\Ml* physiologische Kochsalzlösung 
  (\ce{NaCl} 0,9\PC) 
  {\dotfill} Konzentration 1\,\nicefrac{\Mg}{\Ml}
\end{itemize}

Manche Medikamente liegen als \I{Trockensubstanz} in einer
\I{Trockenstechampulle} vor, welche vor Verwendung erst mit
einem \EE{geeigneten} Lösungsmittel aufgelöst werden muss.
Das erforderliche Lösungmittel ist vom Medikament abhängig.
Bei Verwendung eines falschen Lösungsmittel kann es z.\,B.
zu Ausflockungen, Wirkminderungen o.\,ä. kommen, das
Medikament ist damit verdorben und darf keinesfalls
verwendet werden.

}
{}
{}
{}
{}

\begin{Danger}
  \item Medikamente sind oft auch mit \EE{gleichem Namen} in 
  unterschiedlichen Konzentrationen erhältlich 
  (z.\,B. Dormicum{\texttrademark}, Ketanest{\texttrademark}).
\end{Danger}
\begin{Synopsis}
  \SynopsisItem{Beim Zubereiten bzw. Aufziehen von Medikamenten muss 
  man sehr genau auf die Wirkstoffmenge, das Volumen und die 
  Konzentration achten.}
  \SynopsisItem{Beim Auflösen oder Verdünnen von Medikamenten muss 
  darauf geachtet werden, eine geeignete Flüssigkeit zu verwenden!}
\end{Synopsis}

% \begin{PictureSeriesWide}{}{Bilderserie: Medikamente \SourcePicture{ASBÖ/gjh}}
% \PictureSeriesRowWide{3}
% % {./images/ASBOE-Grassl-gjh-1/Medikamente0395_0197-00441pt}
% % {L-Adrenalin \enquote*{Fresenius\texttrademark} 2\Mg* 20\Ml* Ampulle. 
% % Seitlich ist ein roter Punkt angebracht, welcher die Sollbruchstelle anzeigt.}
% {./images/ASBOE-Grassl-gjh-1/suprarenin-freigestellt-borderadded-00441pt}
% {Suprarenin{\texttrademark}-Ampulle
% (1\Mg* Adrenalin in 1\Ml*, dieses Medikament wird oft vor Gebrauch verdünnt);
% Der Punkt zeigt die Sollbruchstelle an.}
% {./images/ASBOE-Grassl-gjh-1/Medikamente0422_0224-00441pt}
% {Beispiel für eine Trockenstechampulle: 
% Solu-Dacortin{\texttrademark} 250\Mg* Trockenstechampulle: 
% \enquote{Zur Injektion nach Auflösen in 5\Mg* Aqua ad injectionem}}
% {./images/ASBOE-Grassl-gjh-1/Medikamente0405_0207-00441pt}
% {Eine Infusionsflasche zur Kurzinfusion mit Glukose 33\PC}
% {}
% {}
% \end{PictureSeriesWide}

% TODO Grassl Foto: Bilderserie Medikamente

\subsection{\HeadingKeyword{Verabreichungsarten}}
% \TextSlide{\TxBeschreibung}{
Es gibt verschiedene Arten Medikamente zu verabreichen. Dies
kann entweder \II{invasiv} (verletzend) oder
\II{nicht-invasiv} sein. 
% }{%
% \begin{itemize}
%     \item Invasiv
%     \item Nicht-invasiv
% \end{itemize}
% }
% {}{}{}

\TextSlide{Invasiv}{%
\index{Injektion!intravenöse}%
\index{Injektion!intramuskuläre}%
\index{Injektion!subkutane}%
\index{Injektion!intraossär}%
\index{intravenös}%
\index{intramuskulär}%
\index{intraossär}%
Bei der invasiven Verabreichung wird das Medikament mittels
einer Kanüle in das jeweilige Gewebe bzw. Gefäß dem
Patienten injiziert. Diese Verabreichungsarten sind
i.\,d.\,R.
den Ärzten, besonders dafür geschulten
Sanitätern\VARIANT{REG-AUT}{ (NKV)}, diplomiertem
Krankenpflegepersonal oder Hebammen vorbehalten.
Die gängigen Injektionsarten sind: 

\begin{itemize}
  \item \EE{Intravenöse Injektion} (\II{i.\,v.}): Das
  Medikament wird mittels Kanüle, Venenverweilkanüle
  o.\,ä. direkt in eine \E{Vene} gespritzt. 
  \item \EE{Intramuskuläre Injektion} (\II{i.\,m.}): Das
  Medikament wird in den \E{Muskel} verabreicht.
  \item \EE{Subkutane Injektion} (\II{s.\,c.}): Die
  Substanz wird in die Fettschicht  der Unterhaut (Subkutis)
  eingespritzt.
  \item Eine Besonderheit ist der \EE{intraossäre Zugang}
  (\II{i.\,o.}), bei dem die Substanz mittels einer
  speziellen Kanüle in die Knochenmarkhöhle verabreicht wird.
\end{itemize}



}{
\begin{itemize}
    \item Intravenös (i.\,v.)
    \item Intramuskulär (i.\,m.)
    \item Subkutan (s.\,c.)
    \item Intraossär (i.\,o.)
\end{itemize}
}{

}{}{}

\TextSlide{Nicht invasiv}{
Bei der nicht invasiven Verabreichung wird das Medikament
über eine Schleimhaut oder direkt über die Haut aufgenommen.
Für die \II{orale} Aufnahme muss das Arzneimittel geschluckt
werden. Bei der \I[sublingual]{sublingualen} Methode kommt
die Substanz unter die Zunge des Patienten. Bei der
\I[perkutan]{perkutanen} Verabreichung wird das Medikament
auf die Haut aufgetragen und aufgenommen.
}{\begin{itemize}
    \item Oral
    \item Sublingual
    \item Perkutan
\end{itemize}}{

}{}



% \clearpage
\subsection{\HeadingKeyword{Aufziehen} eines Medikamentes aus einer 
Ampulle}

\Text{\TxBeschreibung}{Glasmpullen lassen sich
unterschiedlich öffnen:
Ampullen mit einem \EE{roten Punkt} haben eine
\E{Sollbruchstelle}, der Daumen muss beim Aufbrechen auf dem
roten Punkt liegen. Ampulle mit einem \enquote{\EE{Halsring}}
können nach allen Richtungen aufgebrochen werden.
Beim Arbeiten muss besonderes Augenmerk auf eine \EE{sterile
Arbeitsweise} gelegt werden! Es kann sonst zu Infektionen
kommen.
}{}{}{}{}

\begin{ProcedureStepsNoHeading}
\Text{Material}{
\begin{compactitem}
  \item Spritze, passende Größe
  \item Kanüle(n) 
  \item Ampulle mit Medikament
  \item Trockene Tupfer
  \item Nadelabwurfbehälter
  \item Ggfs. Verschlussstoppel, Permanentstift
\end{compactitem}
}{}{}{}{}
\end{ProcedureStepsNoHeading}

% \begin{multicols}{2}
% \Text{Vorgehen}{
\begin{ProcedureStepsWide}{Vorgehen}
\begin{enumerate}\CorrectionItemizeBox
  \item Verschlossene \StepHeading{Ampulle vorzeigen} bevor sie geöffnet wird.
  \item \StepHeading{Kontrolle} Ablaufdatum und Inhalt: Verfärbung,
  Ausflockung?
  \item \StepHeading{Öffnen} der Ampulle:
  \begin{itemize}
    \item  Glasampulle:

    Ampullenspitze beklopfen, darauf achten, dass sich
    keine Flüssigkeit im Ampullenkopf mehr befindet

    Öffnen der Medikamenten-Ampulle mit Tupfer, roter Punkt
weist zum Körper

    \item Plastikampulle mit
    \enquote*{Dreh-und-Drink}-Verschluss:
    Verschlussknopf abdrehen.

  \end{itemize}
  \item Ampulle abstellen.
  \item \StepHeading{Spritzenverpackung öffnen}: aufschälen, Spritze in
  Verpackung lassen.
  \item \StepHeading{Kanülenverpackung öffnen}: aufschälen 
  \item \StepHeading{Spritze auf Kanüle}
  aufstecken.
  \item \StepHeading{Kanülenabdeckung entfernen} (nicht abdrehen,
  sondern abziehen!)
  \item \StepHeading{Kanüle in Ampulle} einführen

  Beim Einführen der Kanüle in die Ampulle den
    Ampullenrand nicht berühren.
  \item \StepHeading{Aufziehen} des
  Medikamentes.
  \item Verabreichung:
  \ListBreak{enumerate}{Variante \enquote{Verabreichung über
  Venenzugang}:}
    \item Kanüle \EE{händisch abziehen} und in den
    Nadelabwurfbehälter verwerfen -- \Ca{Keinesfalls Kanüle am
    Nadelabwurfbehälter abstreifen (Kontamination)!}
    \item \StepHeading{Entlüften}
    \item Wird das Medikament nicht sofort verwendet, Spritze mit
    Stoppel verschließen und Ampulle ankleben
    \item Spritze übergeben, Ampulle zeigen
\ListBreak{enumerate}{Variante \enquote{zuspritzen in eine
Infusionsflasche}:}
    \item Gummimembran desinfizieren.
    \item Kanüle nicht bis zum Schaft einstechen.
    \item Medikament zuspritzen.
    \item Beschriften der Flasche.
\ListBreak{enumerate}{Variante \enquote{direktes Spritzen}}
    \item \StepHeading{Kanülengröße} erfragen
    \item \StepHeading{Aufziehkanüle} \Ca{händisch} entfernen
    \item \StepHeading{Entlüften}
    \item \StepHeading{Spritzkanüle} aufstecken
    \item Spritze mit Schutzkappe über Kanüle reichen, Ampulle zeigen
\end{enumerate}
\end{ProcedureStepsWide}
% }{}{}{}{}


\begin{PictureSeriesWide}{}{Bilderserie: \EEE{Medikament aus
Ampulle aufziehen} \SourcePicture{Motal}}
\PictureSeriesRowWide{4}
{./images/motal-aass/medi-kontrolle_4807-00800.jpg}
{Kontrolle}
{./images/motal-aass/medi-aufziehen_4815-00800.jpg}
{Aufziehen mit Kanüle}
{./images/motal-aass/medi-kontrollieren_4824-00800.jpg}
{Übergabe}
{./images/motal-aass/medi-verabreichen_4826_CLOSEUP-00800.jpg}
{Verabreichung über Venflon durch befugtes Personal}
\end{PictureSeriesWide}



% \clearpage
\subsection{Vorbereitung einer \HeadingKeyword{Infusion}}
\label{topic:infusion}

\begin{ProcedureStepsNoHeading}
\Text{Material}{%
\Married{
\begin{compactitem}
    \item Infusion (Flasche, Beutel)
    \item Infusionsbesteck
    \item Aufhängevorrichtung
    \item Alkoholhaltiges Desinfektionsmittel
    \item einige saubere Tupfer
    \item Permanentmarker
\end{compactitem}
}{
% \centering\includegraphics[width=0.5\linewidth]{./images/infusion_tropfkammer.jpg}
}
}
{}
{}
{}
{}
\end{ProcedureStepsNoHeading}



\begin{ProcedureStepsWide}[Beim Arbeiten muss besonderes Augenmerk auf eine sterile
Arbeitsweise gelegt werden!]{Vorgehen}

\begin{enumerate}
  \item Ggfs. \StepHeading{Aufhängevorrichtung} auf Flasche
  aufsetzen 
  \item \StepHeading{Abdeckung} entfernen (Membranschutz)
  \item Ggfs. \StepHeading{Wischdesinfektion} der Gummimembran
  
  Alkoholgetränkter Tupfer kann während weiterer Vorbereitung
    auf Membran belassen werden.
  \item \StepHeading{Infusionsbesteck montieren}
  \begin{enumerate}
    \item \StepHeading{Durchflußregler schließen} (Radklemme)
    \item Einstich des Tropfkammerdorns in Infusionsflasche mit
    geschlossenem Lüftungsventil und geschlossener Radklemme
    \item Dorn der Tropfkammer und Gummimembran der
    Infusionsflasche dürfen nicht berührt werden!
  \end{enumerate}
  \item \StepHeading{Infusionsschlauch füllen}
  \begin{enumerate}
    \item Flasche wenden und hochhalten
    \item Tropfkammerluftventil öffnen
    \item Tropfkammer durch Zu\-sam\-m\-en\-drück\-en bis zur Hälfte füllen
    \item Konusabdeckung des Infusionsschlauchs am anderen
    Ende muss bei der Belüftung nicht entfernt werden
    \item Infusionsflasche und Ende des Infusionsschlauchs
    mit nicht dominanter Hand halten
    \item Radklemme öffnen
    \item Infusionsschlauch entlüften
    \item Radklemme schließen
  \end{enumerate}
\end{enumerate}
\end{ProcedureStepsWide}

\begin{PictureSeriesWide}{}{Bilderserie: Material für Infusion}
\PictureSeriesRowWide
{2}% Anzahl Spalten
{./images/pallinger-christoph-aass/Infusionsbesteck_33070-00943pt}% Pfad Bild 1
{Infusionsbesteck \SourcePicture{Christoph Pallinger}}% Text Bild 1
{./images/pallinger-christoph-aass/Infusionsbeutel_33016-00943pt}% Pfad Bild 2
{Beispiele für Infusionsbeutel und -flaschen \SourcePicture{Christoph Pallinger}}% Text Bild 2
{}% Pfad Bild 3
{}% Text Bild 3
{}% Pfad Bild 4
{}% Text Bild 4
%
\PictureSeriesRowWide
{4}% Anzahl Spalten
{./images/gabriel-sebastian-aass/Infusion-Closeup-Tropfkammer-Live-00441pt}% Pfad Bild 1
{Der Dorn des Infusionsbestecks mit der Tropfkammer wird an der Infusionsflasche oder -beutel befestigt, die Tropfkammer ca. bis zur Hälfte gefüllt.}% Text Bild 1
{./images/gabriel-sebastian-aass/Infusion-Closeup-Stellrad-Live-00441pt}% Pfad Bild 2
{Mit der Radklemme wird die Infusionsgeschwindigkeit reguliert.}% Text Bild 2
{./images/gabriel-sebastian-aass/Infusion-Closeup-3WegHahn-Live-00441pt}% Pfad Bild 3
{Mittels eines 3-Weg-Hahns können weitere Infusionen angeschlossen werden.}% Text Bild 3
{./images/gabriel-sebastian-aass/Infusion-Closeup-AnschlussVenflon-Live-00441pt}% Pfad Bild 4
{Schlussendlich wird die Infusion an einem Zugang zum Patienten angeschlossen, hier an einer peripheren Venenverweilkanüle.}% Text Bild 4
\end{PictureSeriesWide}

\clearpage % TEMP temp clearpage
% \bigskip
\section{\HeadingKeyword{Venenverweilkanüle}: Assistenz}
\label{topic:venflon}

\TextSlide{\TxBeschreibung}{
% (\enquote{Venflon}): Vorbereitung und Assistenz 
Periphere Venenverweilkanülen werden routinemäßig zur
intravenösen Gabe von Medikamenten bzw. Infusionen
verwendet. Sie bestehen aus einem Kunststoffschlauch, der in
der Vene zu liegen kommt, und aus einem außenstehenden Teil
mit einer Konnektionsstelle für Infusionen, sowie oft einem
Konus mit Rückschlagventil, für die schnelle Gabe von
Medikamenten. Im Kunststoffschlauch liegt eine
Führungskanüle (\I{Mandrin}), welche das Durchstechen der
Haut und der Vene zwecks der Katheterisierung ermöglicht.
Die Kanüle wird nach Einführen ins Blutgefäß entfernt.

Gebräuchliche Marken sind z.\,B. \I{Venflon{\texttrademark}}
oder \I{Braunüle{\texttrademark}}. Venenverweilkanülen gibt
es in verschiedenen Durchmessern, sie sind farbcodiert.
Normalerweise werden \EE{grüne} Venenverweilkanülen bei
Erwachsenen verwendet, bzw. rosafarbene bei schlechten
Venenverhältnissen.
}{
\begin{itemize}
  \item Kunststoffschlauch, der in
der Vene zu liegen kommt
\item 
Konnektionsstelle für Infusionen
\item Konus mit Rückschlagventil
\item Führungskanüle (\I{Mandrin})
\item Verschiedene Größen, farbcodiert.
Standard: \EE{grün}, rosa bei schlechten
Venen
\item Gebräuchliche Marken:
Z.\,B. \I{Venflon{\texttrademark}},
\I{Braunüle{\texttrademark}}
\end{itemize}

}{}{}{}

\begin{figure}[H]
\includegraphics[width=\linewidth]{./images/demaw/venflon_zerlegt-textwidth.jpg}

\Captionof{figure}{Eine periphere Venenverweilkanüle
(Venflon\texttrademark), in ihre Einzelbestandteile zerlegt.
\SourcePicture{DEMAW}}
\end{figure}



\begin{table}[H]\TableBaseModification
\Captionof{table}{Verschiedene Größen von peripheren
Venenverweilkanülen. \Source{DEMAW}}

\begin{tabular}{m{1.5619999cm}cccll}
\toprule\addlinespace
\TabHeading{Farbcode}
&

&
\TabHeading{Gauge}
&

&
\TabHeading{Länge}
&
\TabHeading{Durchfluss}
\\
 		& 				&  			&{\O} \MM 		& \MM  		& [\nicefrac{ml}{min}] 
\\\addlinespace\midrule\addlinespace
blau 	& Sehr dünn		& \Z{22} 	&\Z{0,5 x 0,8} 	& \Z{25} 	& \Z{25} 
\\
rosa 	& dünn			& \Z{20} 	&\Z{0,7 x 1,0} 	& \Z{32} 	& \Z{55} 
\\
grün 	& normal		& \Z{18} 	&\Z{0,9 x 1,2} 	& \Z{40} 	& \Z{90} 
\\
weiß 	& dick			& \Z{17} 	&\Z{1,1 x 1,4} 	& \Z{42} 	& \Z{135} 
\\
grau 	& Sehr dick		& \Z{16} 	&\Z{1,3 x 1,7} 	& \Z{42} 	& \Z{170} 
\\
orange 	& Noch dicker	& \Z{14} 	&\Z{1,6 x 2,1} 	& \Z{42}
& \Z{265} \\\addlinespace\bottomrule
\end{tabular}
\end{table}


\begin{ProcedureStepsNoHeading}
% \clearpage
% \begin{multicols}{2}
\Text{Material}{
\begin{compactitem}
  \item Nierentasse
  \item Staubinde
  \item Venflon:
  \begin{compactitem}
    \item Die Größe (= Farbe) ist vom Durchführenden vorher zu erfragen!
    \item Die Flügel werden hinunter geklappt
    \item Die farbige Lasche des Zuspritzventils in rechten Winkel zur
    Venflonachse stellen
  \end{compactitem}
  \item Einmalverschlusskappe (\enquote{roter
  Stoppel})
  \item Fixierungspflaster und evtl. zusätzliches
  Fixationsmaterial
  \item Trockene Tupfer
  \item Desinfektions-Tupfer
  \item Spülung, wahlweise:
   \begin{itemize}
    \item Spritze mit \Range{5}{10}\Ml* Spülflüssigkeit
  (\ce{NaCl} 0,9\PC*), oder
  \item Infusion
  \end{itemize}
  \item 3-Wege-Hahn (optional)
  \item (Faserschreiber)
  \item Nadelabwurfbehälter
\end{compactitem}
% \end{multicols}
}{}{}{}{}
\end{ProcedureStepsNoHeading}

\begin{ProcedureStepsNoHeading}
% \begin{multicols}{2}
\Text{Assistenz}{
\begin{enumerate}
  \item Ggfs. \StepHeading{3-Wege-Hahn spülen}
  \begin{itemize}
    \item Spritze mit Spülflüssigkeit ansetzen und durchspülen 
    \item Alternativ, bei Verwendung mit Infusion: 3-Wege-Hahn an Infusionsbesteck anschließen und mit Infusionslösung durchspülen
  \end{itemize}
  \item \StepHeading{Nadelabwurfbehälter} in Reichweite des  Durchführenden
  bereitstellen, sodass er später den Mandrin abwerfen kann.
  \item \StepHeading{Stauschlauch} zureichen
  \item \Z{Der Durchführende legt nun den Stauschlauch an, staut, und sucht eine passende Vene für die Punktion.}
  \item \StepHeading{Alko-Tupfer zureichen}: Mit geöffneter Verpackung sodass dieser
  entnommen werden kann
  \item \StepHeading{Venenverweilkanüle zureichen}:   
  An der Schutzkappe halten, sodass der
  Durchführende die Verweilkanüle abziehen kann
  \item \Z{Der Durchführende punktiert nun die Vene.}
  \item \StepHeading{Trockene Tupfer} zureichen
  \item Je nach Situation 3-Wege-Hahn, Stoppel und Spülung oder Infusionsschlauch
  zureichen
  \item \StepHeading{Fixierpflaster} zureichen
\end{enumerate}
% \end{multicols}
}{}{}{}{}
\end{ProcedureStepsNoHeading}

\begin{PictureSeries}{}{Bilderserie: \EEE{Periphere Venenverweilkanüle}}
\PictureSeriesRow{2}
{./images/demaw/Skriptumaegf-img35}
{Anlage einer periph. Venenverweilkanüle \SourcePicture{DEMAW/Blacky}}
{./images/demaw/Skriptumaegf-img36}
{Eine fixierte Venenverweilkanüle mit Infusion
\SourcePicture{DEMAW/Blacky}} {./images/demaw/venflon_zerlegt}
{\SourcePicture{DEMAW/Blacky}}
{empty}
{}
{empty}
{}
\end{PictureSeries}

\A{Butterfly gehört erklärt}
% TODO Butterfly


\clearpage
% \bigskip
\section{\HeadingKeyword{Endotracheale Intubation}: Vorbereitung und
Assistenz} 
\label{topic:intubation-assistenz}%

\noindent Bei der endotrachealen Intubation wird ein
Beatmungsschlauch (\I{Tubus}) über den Mund (seltener über die Nase) durch den
Rachen und den Kehlkopf in die Luftröhre eingeführt. Er
ermöglicht die Beatmung des Patienten bei gleichzeitigem
Freihalten des Atemweges und bietet einen Aspirationsschutz.

\Text{Material und Personal}{
\begin{compactitem}
  \item 2 Helfer:
  \begin{compactitem}
    \item \ActorAction{ActorA}{A}{Assistenz}
    \item \ActorAction{ActorN}{N}{Durchführender}
  \end{compactitem}
\end{compactitem}
}

\begin{ProcedureStepsWide}[Kontrolle vor Beginn
der Intubation nicht vergessen!]{Material}
\begin{enumerate}
    \item Beatmungsbeutel inkl. passender Beatmungsmaske,
    Bakterienfilter und Sauerstoffreservoir
    \item \ce{O2}-Berieselungseinheit und \ce{O2}-Line
    \item Evtl. Beatmungsgerät
    \item Absaugeinheit: \par\Ca{Keine Intubation ohne
    Absaugung!}
    \item Laryngoskop, bestehend aus:
    \begin{enumerate}
        \item Spatel: gebogen oder gerade; Größen 0\,--\,4
        \item Handgriff mit Batterien
    \end{enumerate}
    \Z{Es gibt Warm- und Kaltlichtlaryngoskope.
    Bei Warmlichtlarynkoskopen befindet sich die Lichtquelle
    am Spatel, bei Kaltlichtlaryngoskpen befindet sich die
    Lichtquelle im Handgriff, das Licht wird über Fiberglas
    in den Spatel geleitet.} \Ca{{\SymbolGefahren}} Daher muss
    darauf geachtet werden, dass die Spatel mit dem Griff
    zusammenpassen! Bei Kaltlichtsystemen gibt es ausserdem
    unterschiedliche Steckverbindungen!
    
    \noindent\begin{minipage}{\linewidth}
\includegraphics[width=\linewidth]{./images/pallinger-christoph-aass/Testlunge_32902-AASS-0112mm.jpg}

\Captionof{figure}{Laryngoskopgriff und verschiedene Spatel
\SourcePicture{Ch. Pallinger}}
\end{minipage}
    \item (Endotracheal-)Tubus; verschiedene Größen
    \item Führungsdraht (\T{Mandrin})
    \item Silikonspray
    \item Evtl. Magillzange
    \item Mit Luft gefüllte Spritze zum Cuffen (je
    nach Tubus, i.\,d.\,R. 10\Ml*)
    \item Stethoskop
    \item Beißschutz: Beißkeil oder Guedel-Tubus
    \item Fixationsmaterial (Breite Heftpflaster, Mullbinde)
\end{enumerate}


\end{ProcedureStepsWide}





\PictureWide
{Zubehör für die endotracheale Intubation}%1 Titel
{Absaugeinheit
mit Absaugkatheter, Magill-Zange, Silikonspray, Tubus,
Führungsdraht, Blockerspritze (10\Ml*), verschiedene Spatel,
Laryngoskop, Guedel-Tuben, Beißkeil, Beatmungsbeutel mit
Bakterienfilter, Maske und \ce{O2}-Line, Fixationsmaterial,
Stethoskop}%2 Beschreibung
{}%3 Label
{Ch. Pallinger}%4 Quelle
{MfG}%5 Lizenz
{./images/pallinger-christoph-aass/Intubation_32939.jpg}%6 Datei
{}%7 Breitenmodifizierer
{}%8 Float
% vormals \Layout{37}

\clearpage



% \MarginBoxFilledRounded{Vorbereitung}{}

\begin{ProcedureStepsWide}{Vorbereitung}
\begin{enumerate}
  \item \StepHeading{Präoxygenieren}: 
  \begin{compactitem}
     \item Beim spontan atmenden Patienten: 
     \ActorAction{ActorA}{A}{Hochdosierte Sauerstoffgabe über \ce{O2}-Maske}
     \item Beim nicht-spontan-atmenden Patienten: 
     \ActorAction{ActorN}{N}{Je nach Situation evtl. Beatmung mittels Beatmungsbeutel}\ZFootnote{%
     Beim nicht-nüchternen 
     nicht-spontan-atmenden Patienten soll eine Beutelbeatmung nur 
     wenn unbedingt notwendig 
     durchgeführt werden. 
     %
     Bei der Reanimation entfällt die Präoxygenierung.} %ZFootnote
  \end{compactitem}
  
  \item \ActorAction{ActorA}{A}{\StepHeading{Erfragen}:}
  \begin{compactitem}
    \item Spatelgröße 
    \item Tubusgröße
  \end{compactitem}
  %
  \item \ActorAction{ActorA}{A}{\StepHeading{Funktionskontrolle}: 
  Laryngoskop 
  (Spatel und
  Griff) 
  zusammenbauen, 
  Funktionskontrolle (Licht)
  
  Ebenso müssen alle möglicherweise benötigten Spatel
  (zumindest $\pm$ 1 Größe und die entsprechenden
  gebogenen bzw. geraden Spatel) mit dem
  vorhandenen Laryngoskopgriff geprüft werden (Lichtquelle
  bei Warmlichtgeräten, passende Verbindung bei Kaltlichtgeräten)}
  \item \ActorAction{ActorA}{A}{Tubusverpackung aufschälen, Tubus aber in der Verpackung belassen}
  \item \ActorAction{ActorA}{A}{\StepHeading{Cuff-Dichtigkeit} prüfen: test\-weise aufblasen}
  \item \ActorAction{ActorA}{A}{\StepHeading{Führungsdraht} (Mandrin) mit Silikonspray
  einsprühen\ACh{HART}{0.8.0}{Mandrin auch mit Silikon einsprühen} und
  einführen. 
  Die Spitze muss bis zum Tubusende vorgeschoben vorgeschoben
  werden, darf aber nicht herausragen.}
  \item \ActorAction{ActorA}{A}{\StepHeading{Cuff} mit Silikonspray einsprühen}
  \item \ActorAction{ActorA}{A}{\StepHeading{Beatmungsbeutel} zusammenbauen und an
  \ce{O2}-Berieselungseinheit anschließen}
  \item \ActorAction{ActorA}{A}{\StepHeading{Stethoskop}: Funktion prüfen}
  \item \ActorAction{ActorA}{A}{\StepHeading{Absaugungebreitschaft} herstellen:}
  \begin{enumerate}
    \item Passenden Absaugkatheter auspacken, anstecken und sauber ablegen
    \item Gerät einschalten
    \item Funktionskontrolle: Saugstärke?
    Batteriewarnung?
    \item Absaugeinheit und Absaugkatheter muss in Griff\-weite des Assistenten
    positioniert werden!
  \end{enumerate}
  \item \ActorAction{ActorA}{A}{Restliches Material bereitlegen (z.\,B. in Nierentasse)}
  \ListBreakHard[Für nicht-ärztliches
  Personal]{enumerate}{Assistenz}
  \item 
  \Actor{ActorA}{A}\Actor{ActorN}{N} \StepHeading{Material-Vollständigkeit} prüfen
  \item \ActorAction{ActorA}{A}{\StepHeading{Sauerstoffberieselung} einschalten und auf
  15{\LPerMin*} einstellen, Reservoir des Beatmungsbeutels
  füllen}
  \item \ActorAction{ActorA}{A}{Durchführendem melden, dass alle Materialien bereit sind}
  \item 
  \ProcedureStepMargin{Medikamentengabe (Narkoseinleitung)}
  \ActorAction{ActorN}{N}{Es wird nun von einem Arzt die Narkoseeinleitung durchgeführt
  (entfällt bei der Reanimation).
  Dazu werden Medikamente
  verabreicht, 
  die die Intubation ermöglichen.} 
  
  Von jetzt an besteht ein sehr hohes Risiko dass es zu \E{Erbrechen} kommen kann,
  ggfs. muss sofort abgesaugt werden!
  %
  Der Wirkeintritt der Medikamente wird abgewartet.
  %
%  Der
%  Durchführende wird nun einige Zeit den Patienten mittels
%  Beutel-/Masken-Beatmung beatmen (sofern nicht bereits
%  geschehen).
  \item \ActorAction{ActorA}{A}{\StepHeading{Laryngoskop zureichen}: 
  \marginpar{Laryngoskop zureichen}
  in die \E{linke (!)}
  Hand zureichen\footnote{Es ist egal, ob der Durchführende
  Rechts- oder Linkshänder ist, das Laryngoskop wird
  \E{immer in der linken Hand} gehalten.}}
  \item \ActorAction{ActorA}{A}{Auf Anweisung \emph{\enquote{Kehlkopfdruck}
  (\enquote{Krikoiddruck}}) durchführen:
  auf den Kehlkopf drücken}
  \item \ActorAction{ActorA}{A}{\StepHeading{Tubus zureichen}: 
  in die \E{rechte (!)} Hand}
  \item \ActorAction{ActorA}{A}{Auf Anweisung absaugen}
  \item \ActorAction{ActorN}{N}{Der Tubus wird nun soweit eingeführt, 
  dass die Tubusspitze zwischen die Stimmritzen durchdringt.}
  \item \ActorAction{ActorN}{N}{Wenn die Tubusspitze durch die Stimmritzen durchgetreten ist muss der Mandrin entfernt werden um Verletzungen zu vermeiden. Anweisung: \Speech{\enquote{Führungsdraht entfernen!}}} 
  \item 
  \ProcedureStepMargin{Mandrin herausziehen}
  \ActorAction{ActorA}{A}{Der Mandrin wird herausgezogen, 
  dabei darf keinesfalls am Tubus gezogen werden, 
  da dieser nur knapp in die Luftröhre hineinragt.}
  %
  \item \ActorAction{ActorN}{N}{Der Tubus wird nun bis zur bis zu endgültigen Tiefe eingeführt 
  und vorübergehend manuell fixiert.}
  %
  \item 
  \mbox{}\ProcedureStepMargin{Cuffen}
  \ActorAction{ActorN}{N}{Anweisung: 
  \Speech{\enquote{Cuffen}}}
  %
  \item \ActorAction{ActorA}{A}{\StepHeading{Cuffen}: 
  Mit luftgefüllter Spritze Cuff aufblasen}
  %
  \item \ActorAction{ActorA}{A}{\StepHeading{Stethoskop} 
  dem Durchführenden in die Ohren klemmen}
  \item\ActorAction{ActorA}{A}{Beatmungsmaske vom Beatmungsbeutel trennen}
  \item \ActorAction{ActorA}{A}{\StepHeading{Beatmungsbeutel} an Tubus anschließen und halten}
  \item 
  \ProcedureStepMargin{Lagekontrolle}
  \StepHeading{Lagekontrolle} mittels Auskulatation: 
  \ActorAction{ActorN}{N}{Magengegend und beide Lungenflügel auskultieren,
  dabei jedesmal mindestens einen Atemhub mittels Anweisung 
  \Speech{\enquote{Atemhub}} 
  verabreichen lassen.}
  \ActorAction{ActorA}{A}{Auf Anweisung werden
  mindestens 3 Atemhübe verabreicht}
  \item \ActorAction{ActorN}{N}{Der Durchführende übernimmt den Beatmungsbeutel.}
  \item \ActorAction{ActorA}{A}{Einführen des \StepHeading{Beißkeils} oder des Guedel-Tubus}
  \item \ActorAction{ActorA}{A}{\StepHeading{Fixierung} des Tubus und des Beißschutzes
  mittels Mullbinde, Klebestreifen oder speziellen
  Fixationsmaterialien}
  \item \Z{Ggfs. Beatmungsgerät einstellen und
  anschließen}
\end{enumerate}
\end{ProcedureStepsWide}
% }{}{}{}{}


% \begin{PictureSeriesWide}{}{Bilderserie: Intubation \SourcePicture{Motal}}
% \PictureSeriesRowWide{4}
% {./images/motal-aass/intubation/dsc_4829-AASS-0030mm}
% {Zusammenbauen des Laryngoskops}
% {./images/motal-aass/intubation/dsc_4831-AASS-0030mm}
% {Einsprühen des Tubus \E{in der Verpackung}}
% {./images/motal-aass/intubation/dsc_4834-AASS-0030mm}
% {Präoxygenierung}
% {./images/motal-aass/intubation/dsc_4835-AASS-0030mm}
% {Zureichen des Laryngoskops}
% \PictureSeriesRowWide{4}
% {./images/motal-aass/intubation/dsc_4837-AASS-0030mm}
% {Zureichen des Tubus}
% {./images/motal-aass/intubation/dsc_4839-AASS-0030mm}
% {Einführen des Tubus, Assistenz führt auf Anweisung Kehlkopfdruck aus}
% {./images/motal-aass/intubation/dsc_4841-AASS-0030mm}
% {Cuffen nach Entfernung des Führungsdrahtes}
% {./images/motal-aass/intubation/dsc_4843-AASS-0030mm}
% {Stethoskop zureichen}
% \PictureSeriesRowWide{4}
% {./images/motal-aass/intubation/dsc_4844-AASS-0030mm}
% {Abhören des Patienten}
% {./images/motal-aass/intubation/dsc_4847-AASS-0030mm}
% {Beißkeil (hier: Guedel-Tubus) einführen}
% {./images/motal-aass/intubation/dsc_4849-AASS-0030mm}
% {Endgültige Fixierung des Tubus: Bis dahin muss der Tubus manuell fixiert werden!}
% {empty}
% {}
% \end{PictureSeriesWide}

\begin{PictureSeriesWide}{}{Bilderserie: \EEE{Intubation} \SourcePicture{Michael Motal}}
\PictureSeriesRowWide{3}
{./images/motal-aass/intubation/dsc_4829-AASS-0030mm}
{Zusammenbauen des Laryngoskops}
{./images/motal-aass/intubation/dsc_4831-AASS-0030mm}
{Einsprühen des Tubus \E{in der Verpackung}}
{./images/motal-aass/intubation/dsc_4834-AASS-0030mm}
{Präoxygenierung}
{empty}
{}
\PictureSeriesRowWide{3}
{./images/motal-aass/intubation/dsc_4835-AASS-0030mm}
{Zureichen des Laryngoskops (\EE{linke} Hand!)}
{./images/motal-aass/intubation/dsc_4837-AASS-0030mm}
{Zureichen des Tubus (\EE{rechte} Hand!)}
{./images/motal-aass/intubation/dsc_4839-AASS-0030mm}
{Einführen des Tubus, Assistenz führt auf Anweisung Kehlkopfdruck aus}
{empty}
{}
\PictureSeriesRowWide{3}
{./images/motal-aass/intubation/dsc_4841-AASS-0030mm}
{Cuffen nach Entfernung des Führungsdrahtes}
{./images/motal-aass/intubation/dsc_4843-AASS-0030mm}
{Stethoskop zureichen}
{./images/motal-aass/intubation/dsc_4844-AASS-0030mm}
{Abhören des Patienten}
{empty}
{}
\PictureSeriesRowWide{3}
{./images/motal-aass/intubation/dsc_4847-AASS-0030mm}
{Beißkeil (hier: Guedel-Tubus) einführen}
{./images/motal-aass/intubation/dsc_4849-AASS-0030mm}
{Endgültige Fixierung des Tubus: Bis dahin muss der Tubus manuell fixiert werden!}
{./images/gabriel-sebastian-aass/Monitor-Beatmungseinstellungen-Live-00441pt-4}
{Die Kapnometrie ist wichtig, um die korrekte Lage des Tubus in der Luftröhre zu bestätigen. \SourcePicture{Sebastian Gabriel}}
%{./images/gabriel-sebastian-aass/Monitor-Beatmungseinstellungen-MedumatTransport-Live-00441pt}
%{Anschließend kann die Einstellung eines Beatmungsgerätes erfolgen. \SourcePicture{Sebastian Gabriel}}
{empty}
{}
\end{PictureSeriesWide}

\normalfont

\begin{Literary}
  \fullcite{ArbeitstechnikenAZ1}

  \fullcite{InjektionenLeichtGemacht4}

  \fullcite{FamPropAeGruFert2011}
\end{Literary}



\nocite{ZechaOpGehilfe1}
\ChapterEnd